problem:  
Let \(V^n \subset S^{n+k}\) be a stationary integral \(n\)-varifold whose regular part is a smooth minimal submanifold, and assume that at every singular point \(p \in \mathrm{sing}\, V\) there is a **unique tangent cone**  
\[
C_p = \mathbb{R}_+ \times L_p \subset T_p S^{n+k},
\]
with multiplicity one and smooth stable minimal link \(L_p \subset S^{n+k-1}\).  
Denote by \(\lambda_1(L_p)\) the first nonzero eigenvalue of the Laplace–Beltrami operator on \(L_p\).

**Question:**  
Assume that the tangent cones \(C_{p_j}\) at a sequence of singular points \(p_j \to p_\infty\) converge **as stationary varifolds** (not necessarily smoothly) to the tangent cone \(C_{p_\infty}\).  
Is the map
\[
p \longmapsto \lambda_1(L_p)
\]
necessarily *upper semicontinuous* on \(\mathrm{sing}\, V\)?  
That is, does one have
\[
\limsup_{j\to\infty} \lambda_1(L_{p_j}) \leq \lambda_1(L_{p_\infty})?
\]

More broadly, does the positive curvature of \(S^{n+k}\) (or the stability of the links \(L_p\)) enforce such analytic rigidity on the spectral data of links—so that, even under weak varifold convergence of cones, the first Laplacian eigenvalue cannot increase?

If not, can one construct a sequence of stationary minimal varifolds in \(S^{n+k}\) with unique tangent cones whose links \(L_{p_j}\) converge as varifolds to \(L_{p_\infty}\) but  
\[
\limsup_{j\to\infty} \lambda_1(L_{p_j}) > \lambda_1(L_{p_\infty})?
\]
That is, can spectral discontinuity occur under weak geometric convergence within the positively curved ambient sphere?

Why is it a "good" problem:  
This version isolates the precise analytic subtleties of spectral behavior under **varifold (weak) convergence** of tangent cones, avoiding the false equivalence with smooth convergence. It poses a concrete open question at the interface of *geometric measure theory* (weak convergence of stationary cones) and *spectral geometry* (stability of Laplacian eigenvalues).  

A positive answer would reveal a new rigidity principle, suggesting that the ambient positive curvature or spectral stability of minimal links extends even to the distributional level of tangent cone limits.  
A counterexample would expose hidden analytic flexibility in minimal singularities, highlighting the gap between geometric and spectral convergence.  

The problem is mathematically well-defined (under the unique-tangent-cone assumption), conceptually clear, and free of routine technicalities—it distills a simple, natural invariant (the first eigenvalue of the cone’s link) whose semicontinuity under weak geometric convergence is unknown. It thus exemplifies the “narrow entrance, profound depth” character of a genuinely good research problem, bridging spectral analysis, curvature, and the fine structure of singularities in minimal geometry.